% =======================================================
% 8. PREVENTIVO A FINIRE (PaF)
% =======================================================
\section{Preventivo a Finire (PaF)}
\label{sec:paf}

Il Preventivo a Finire viene aggiornato al termine di ogni sprint a valle del ciclo di controllo del Piano di Progetto. L'aggiornamento tiene conto del consuntivo di periodo, della Sprint Retrospective\textsubscript{G}, delle misure correttive individuate, dell'aggiornamento dei rischi e della revisione della pianificazione futura.

Con la conclusione del progetto allo Sprint 21 (Sezione~\ref{sec:conclusione}), a seguito della validazione formale positiva dell'MVP da parte della proponente (\link{https://synergo-unit-unipd.github.io/Second-Brain/docs/PB/doc_gestionali/verb_esterni/verb_est_2026_09_01/verb_est_2026_09_01_firmato.pdf}{Verbale Esterno N. 10, 2026-09-01}), questa sezione non fornisce più una stima delle risorse necessarie a completare attività future, non essendo pianificati ulteriori sprint di sviluppo, ma il \textbf{consuntivo finale} del progetto: il costo complessivo effettivamente sostenuto a confronto con il budget preventivato in fase di candidatura, e il trattamento delle ore e del costo non consumati.

\subsection{Analisi scostamenti}

Il monitoraggio delle ore residue per ruolo, aggiornato al termine dello
Sprint 21 evidenzia una distribuzione non uniforme degli scostamenti rispetto
alla preventivazione iniziale:

\begin{table}[H]
\centering
\renewcommand{\arraystretch}{1.3}
\begin{tabular}{l r r}
\toprule
\textbf{Ruolo} & \textbf{Ore non consumate} & \textbf{Risparmio} \\
\midrule
Responsabile   & 18{,}58 & 557{,}40~\euro \\
Amministratore & \color{red}{-55{,}00} & \color{red}{-1.100{,}00~\euro} \\
Analista       & 54{,}67 & 1.366{,}75~\euro \\
Progettista    & 88{,}67 & 2.216{,}75~\euro \\
Programmatore  & 39{,}00 & 585{,}00~\euro \\
Verificatore   & 45{,}00 & 675{,}00~\euro \\
\midrule
\textbf{Totale} & \textbf{190{,}92} & \textbf{4.300{,}90~\euro} \\
\bottomrule
\end{tabular}
\caption{Ore e budget non consumati per ruolo a fine progetto (dopo lo Sprint 21)}
\end{table}

Il ruolo di Amministratore\textsubscript{G} resta l'unico a chiudere il progetto con un
disavanzo di ore (-55,00 ore, pari a -1.100,00~\euro), in peggioramento
rispetto alle -46,67 ore rilevate a fine Sprint 20, per effetto del
consuntivo dello Sprint 21 stesso. Questo scostamento,
già rilevato in fase RTB\textsubscript{G} (Sez. RR/RO6 dell'Analisi dei
Rischi), resta di origine prevalentemente RTB e non è stato riassorbito
durante la PB; come discusso in Sezione~\ref{sec:scostamento-economico}, il
suo impatto economico è comunque ampiamente compensato dal risparmio
accumulato sugli altri ruoli e non ha mai messo a rischio il rispetto del
budget complessivo.

Il ruolo di Progettista chiude il progetto con il maggiore avanzo di ore
(88,67 ore, 2.216,75~\euro), seguito da Analista (54,67 ore, 1.366,75~\euro)
e Programmatore (39,00 ore, 585,00~\euro). Responsabile e Verificatore
chiudono rispettivamente a 18,58 ore (557,40~\euro) e 45,00 ore
(675,00~\euro).

Il saldo complessivo a fine progetto (190,92 ore, pari a 4.300,90~\euro)
resta ampiamente positivo: il disavanzo di Amministratore viene compensato
dagli avanzi degli altri ruoli, in particolare Progettista e Analista, e il
progetto si chiude senza alcun superamento del budget economico complessivo
preventivato in fase di candidatura.

\subsection{Analisi dello scostamento economico}
\label{sec:scostamento-economico}

Il forte disallineamento tra lo scostamento in ore e lo scostamento in
valore economico non è casuale, ma deriva da un errore di allocazione
commesso in fase di preventivazione, non da un problema
di esecuzione durante la PB.

I ruoli di Analista e Progettista, entrambi con tariffa oraria di 25~\euro/h,
sono stati preventivati con un numero di ore superiore a quanto
effettivamente necessario per le attività assegnate: a fine progetto
presentano un avanzo rispettivamente di 54,67 ore (1.366,75~\euro)
e 88,67 ore (2.216,75~\euro), il più consistente tra tutti i ruoli.
Il ruolo di Amministratore, con tariffa oraria inferiore (20~\euro/h), è
stato invece preventivato con ore insufficienti rispetto al carico reale di
gestione documentale e strumenti di progetto, risultando nell'unico disavanzo
(-55,00 ore, -1.100,00~\euro).

Poiché il costo economico di uno scostamento dipende dal prodotto tra ore e
tariffa oraria, lo sforamento in ore dell'Amministratore pesa relativamente
poco in termini monetari proprio a causa della sua tariffa più bassa, mentre
il risparmio accumulato sui ruoli di Analista e Progettista, a tariffa più
alta, pesa proporzionalmente di più a parità di ore. Il risultato netto è
che il gruppo, pur avendo speso più ore di Amministratore di quelle
preventivate, ha comunque conseguito un risparmio economico complessivo:
il surplus di ore preventivate su Analista e Progettista compensa
ampiamente, in valore, il deficit di ore su Amministratore.

Questa analisi conferma che il vincolo economico complessivo non è a
rischio nonostante lo scostamento negativo su un singolo ruolo, e che la
causa del disallineamento è da ricondursi a una calibrazione non ottimale
delle ore per ruolo in fase di preventivazione iniziale, piuttosto che a un
sovra- o sotto-impiego di risorse durante la PB. L'ulteriore peggioramento
del deficit di Amministratore durante lo Sprint 21 (-8,33 ore aggiuntive)
è stato assorbito senza compromettere il budget complessivo, grazie ai
consistenti surplus degli altri ruoli.

\subsection{Ripianificazione}

A seguito degli scostamenti descritti, il gruppo ha ribilanciato nel corso
della PB\textsubscript{G} la distribuzione delle ore tra ruoli,
mantenendo invariato il vincolo economico complessivo definito nel preventivo
iniziale. In particolare:
\begin{itemize}
    \item il debito di ore di Amministratore ereditato dalla RTB\textsubscript{G}
    non è stato ricostituito durante la PB. Il delta netto complessivo su Sprint 12 -- Sprint 21
    rimane comunque contenuto (-0,18 ore), confermando l'efficacia della
    pianificazione delle attività amministrative in sede di Sprint
    Planning\textsubscript{G};
    \item le ore residue del ruolo di Analista, resesi disponibili dal
    consolidamento dell'Analisi dei Requisiti dopo la validazione RTB, sono
    state impiegate nella pianificazione delle attività rimanenti, con
    particolare attenzione alla verifica e all'aggiornamento della
    documentazione tecnica per la PB (Specifica Tecnica, Manuale Utente);
    \item i lievi sovraccarichi di Progettista e Responsabile registrati
    negli Sprint 20 e 21, coerenti con l'intensità delle attività di
    aggiornamento della documentazione tecnica in vista della consegna PB,
    sono stati assorbiti dal residuo positivo accumulato su questi ruoli
    (rispettivamente 88,67 e 18,58 ore residue a fine Sprint 21) senza
    necessità di ulteriori interventi correttivi.
\end{itemize}

Il progetto si conclude allo Sprint 21 (01/09/2026), con la validazione
formale positiva dell'MVP da parte della proponente (Sezione~\ref{sec:conclusione}):
non sono pertanto pianificati ulteriori sprint di sviluppo. 

Le 190,92 ore
(4.300,90~\euro) risultanti dal bilancio finale per ruolo non
rappresentano quindi risorse necessarie a completare attività future, ma il
\textbf{risparmio complessivo} rispetto al budget preventivato in fase di
candidatura: il progetto si chiude nettamente entro il vincolo economico
massimo accettato, senza necessità di alcuna ripianificazione strutturale.

Come discusso in Sezione~\ref{sec:scostamento-economico}, tale risparmio non
deriva da una sotto-esecuzione delle attività pianificate, né da una
riduzione di scope, ma da un preventivo iniziale non ottimamente calibrato:
in particolare da una sovrastima sistematica delle ore necessarie ai ruoli
di Analista e Progettista, solo in parte controbilanciata da una
sottostima delle ore di Amministratore. Il gruppo ritiene comunque corretto
mantenere tracciato lo scostamento per ruolo, anziché limitarsi a riportare
il solo saldo economico complessivo, per documentare in modo trasparente
l'origine dell'errore di preventivazione a beneficio di gruppi futuri.

Le azioni correttive adottate durante la PB, confermate efficaci dai dati di
Sprint 12 -- Sprint 21, sono:
\begin{itemize}
    \item maggiore attenzione alla pianificazione delle attività amministrative e
    di gestione durante lo Sprint Planning, evitando che vengano considerate
    attività accessorie o non preventivate;
    \item monitoraggio più frequente delle ore effettive tramite gli strumenti di
    consuntivazione, con verifica del residuo per ruolo al termine di ogni sprint;
    \item mantenimento del budget complessivo iniziale come vincolo massimo non
    superabile, anche in presenza di scostamenti localizzati su singoli ruoli.
\end{itemize}

Tali azioni risultano coerenti con le strategie di mitigazione definite nel
Piano di Progetto e hanno permesso di stabilizzare lo scostamento di
Amministratore durante la PB, pur senza riassorbire il debito ereditato dalla
RTB, preservando in ogni caso il rispetto delle scadenze e del budget
complessivo.

\subsection{Quadro economico finale}

\begin{table}[H]
\centering
\renewcommand{\arraystretch}{1.3}
\begin{tabular}{l r}
\toprule
\textbf{Voce} & \textbf{Importo} \\
\midrule
Budget iniziale preventivato   & 11.665{,}00~\euro \\
Costo totale effettivamente sostenuto & 7.364{,}16~\euro \\
\midrule
\textbf{Risparmio finale}      & \textbf{4.300{,}90~\euro} \\
\bottomrule
\end{tabular}
\caption{Quadro economico finale del progetto, dopo lo Sprint 21}
\end{table}

Al termine dello Sprint 21, ovvero lo sprint finale, risultano sostenuti 7.364{,}16~\euro, pari a circa il 63,1\% del budget complessivo preventivato.
Questo quantitativo di budget rimanente è dato dal fatto che, all'inizio del progetto, il gruppo ha deciso di concentrare le ore sui 
ruoli di analista e progettista. Nel corso del progetto il gruppo ha compreso che alcune task che sembravano appartenere a questi ruoli erano 
più adatte al ruolo di amministratore il quale infatti risulta con un residuo negativo di -55 ore.
Il preventivo quindi risulta avere un valore stimato molto più alto in quanto sia Analista che progettista sono ruoli dal costo orario di 25 euro 
contro i 20 dell'amministratore. Il budget rimanente quindi è più alto di quanto sarebbe stato se la suddivisione dei ruoli fosse stata più accurata.
Tale risparmio, per le ragioni discusse, non rappresenta capacità
 inutilizzata rispetto ad attività pianificate e non svolte, ma l'esito di una calibrazione non ottimale del 
 preventivo iniziale tra i ruoli di Analista/Progettista (sovrastimati) e Amministratore (sottostimato).

\subsection{Rettifica a seguito della revisione del metodo di rendicontazione}
\label{sec:paf-rettifica}

Le tabelle precedenti riportano i valori così come consuntivati al termine
dello Sprint 21. A seguito della revisione del metodo di rendicontazione
condotta dopo il colloquio di revisione del 2026-09-15
(Appendice~\ref{sec:appendice-rendicontazione}), sono state individuate
162{,}83 ore di lavoro effettivamente svolto ma non rendicontato nel corso
del progetto. Questa sezione riporta l'effetto di tale rettifica sul
residuo per ruolo, senza modificare i valori originari riportati sopra, che
restano quelli consuntivati e validati al momento della revisione.

\begin{table}[H]
\centering
\renewcommand{\arraystretch}{1.25}
\begin{tabular}{|l|c|c|}
\hline
\textbf{Ruolo} & \textbf{Ore residue} & \textbf{Ore residue} \\
 & \textbf{pre-rettifica (Sprint 21)} & \textbf{post-rettifica} \\
\hline
Responsabile   & 18{,}58 & 18{,}58 \\
Amministratore & -55{,}00 & -89{,}50 \\
Analista       & 54{,}67 & 33{,}67 \\
Progettista    & 88{,}67 & 75{,}42 \\
Programmatore  & 39{,}00 & 26{,}00 \\
Verificatore   & 45{,}00 & -15{,}42 \\
\hline
\textbf{Totale} & \textbf{190{,}92} & \textbf{48{,}75} \\
\hline
\end{tabular}
\caption{Ore residue per ruolo, pre- e post-rettifica}
\label{tab:paf-rettifica-ore}
\end{table}

\begin{table}[H]
\centering
\renewcommand{\arraystretch}{1.25}
\begin{tabular}{lr}
\toprule
\textbf{Voce} & \textbf{Importo} \\
\midrule
Budget iniziale preventivato & 11.665{,}00~\euro \\
Risparmio pre-rettifica (Sezione~\ref{sec:paf}) & 4.300{,}90~\euro \\
\textbf{Risparmio post-rettifica} & \textbf{1.653{,}33~\euro} \\
\bottomrule
\end{tabular}
\caption{Quadro economico, pre- e post-rettifica}
\label{tab:paf-rettifica-costo}
\end{table}

Il risparmio post-rettifica (1.653{,}33~\euro, pari a 48{,}75 ore) resta
\textbf{positivo}: il budget complessivo preventivato \textbf{non viene
superato}. Rispetto al dato pre-rettifica, il margine si riduce da
4.300{,}90~\euro a 1.653{,}33~\euro, quasi interamente per effetto del
ruolo di Amministratore (da -55{,}00 a -89{,}50 ore) e del ruolo di
Verificatore (da 45{,}00 a -15{,}42 ore), coerentemente con le due voci di
maggior peso individuate in Appendice~\ref{sec:appendice-rendicontazione}
(rispettivamente 36{,}50 e 57{,}58 ore di lavoro non rendicontato). Il
progetto si chiude quindi entro il vincolo economico di candidatura anche
tenendo conto della rettifica.